\documentclass[10pt,a4paper]{amsart}
\usepackage[margin=19mm]{geometry}
\usepackage{amsmath,amssymb,mathtools}
\usepackage[scr=rsfs,cal=euler]{mathalfa}
\usepackage{xfrac}
\usepackage[unicode,hidelinks,bookmarks=false]{hyperref}
\newcommand{\ie}{i.e.\ }
\newcommand{\cf}{cf.\ }
\newcommand{\resp}{respectively\ }
\newcommand{\Subsection}{\subsection}
\providecommand{\nobreakdash}{\nobreak\hbox{-}}
\setlength{\emergencystretch}{2em}
\multlinegap0pt
\usepackage{bm}
\usepackage{bbm}
\usepackage{dsfont}
\usepackage{xspace}
\usepackage{cancel}
\usepackage{mathtools}
\usepackage{amsthm}
\makeatletter
\@ifundefined{equation}{\newtheorem{equation}{Equation}}{}
\@ifundefined{lemma}{\newtheorem{lemma}[equation]{Lemma}}{}
\makeatother
\makeatletter
\expandafter\ifx\csname AMS\endcsname\relax
\newenvironment{AMS}{}{}
\fi
\expandafter\ifx\csname myproof\endcsname\relax
\newenvironment{myproof}[1] {\par\noindent\textit{Proof of {#1}:}}{\hfill$\square$}
\fi
\newcommand{\seq}[1]{{\left\langle \, {#1} \,\right\rangle}}
\makeatother
\title[On edge-direction and compact edge-end spaces]{On edge-direction and compact edge-end spaces}
\author{Gustavo Boska, Matheus Duzi, Paulo Magalh\~aes J\'unior}
\date{Source: arXiv:2503.19088v2 (CC BY 4.0)}
\begin{document}
\maketitle

\noindent\emph{Source and licence.} On edge-direction and compact edge-end spaces. Version arXiv:2503.19088v2 (CC BY 4.0), distributed under the Creative Commons Attribution 4.0 International licence, as recorded in the arXiv licence metadata for this exact version and shown on the abstract page https://arxiv.org/abs/2503.19088v2. Full rights notice: licenses/arxiv-2503.19088-CC-BY-4.0.txt.\par\smallskip

\noindent\textit{Editorial scope.} Counted unit: the Lemma whose source label is \texttt{LEMMA\_TimidEdgeSeparation} in the article (source0.tex lines 1173--1175, source label \texttt{LEMMA\_TimidEdgeSeparation}), delivered together with the article's own complete original proof of it (source0.tex lines 1176--1217, one \texttt{proof} environment of 42 lines). No mathematics is written, repaired or re-derived: statement and proof are the article's own text line for line. Required context retained uncounted: LEMMA\_TimidEqvClassQuot (lines 1155--1157). Every definition, display and label of those lines is retained unchanged. Omitted from this package and not redistributed: 1--1154, 1158--1172, 1218--1493. Nothing inside a retained excerpt is omitted or altered: every retained body is delivered exactly as the article's lines are written, so no omission receipt of any kind accompanies this package. Citations: the retained excerpts carry no citation command at all, so no bibliography entry of the article is transcribed and no reference is redistributed. Reference adaptation: none was needed and none was made; every reference inside the delivered unit resolves inside it, so no unresolved reference or citation is produced. Printed slips kept literal: no printed word, symbol, hypothesis or number of the counted statement or of its proof was altered. Layout and class: the article's own class is \texttt{article} with packages including geometry, amsfonts, amsmath, amscd, amssymb, stmaryrd, amsthm, centernot, cmap, fontenc, inputenc, lmodern, graphicx, caption; the wrapper is the lane's pinned \texttt{amsart} editorial wrapper at 10pt on A4, which restates the theorem-like environments the retained text needs and adds the article's own macro definitions that the retained text needs (\texttt{\textbackslash seq}), with extra wrapper packages (bm, bbm, dsfont, xspace, cancel, mathtools). The delivered numbering is the unit's own. Assets: no retained span contains a figure environment, a graphics-inclusion command, a PDF-page inclusion, a table or a TikZ picture; no image file is copied into this package, the package declares no asset dependency and no image file is redistributed with it. Licence: version arXiv:2503.19088v2 of this article is distributed under the Creative Commons Attribution 4.0 International licence, as recorded in the arXiv licence metadata for this exact version and shown on the abstract page https://arxiv.org/abs/2503.19088v2. A full rights notice is delivered at licenses/arxiv-2503.19088-CC-BY-4.0.txt. Characters: every retained excerpt is pure ASCII, so the delivered engine, XeLaTeX with its default Latin Modern text font, reports no missing character.
\begin{lemma}\label{LEMMA_TimidEqvClassQuot}
    Suppose $u\nsim_E v$. Then $u\sim_{E_\mathrm{t}} v$ in $G$ if, and only if, $[u]_E\sim_{E_\mathrm{t}} v$ in $G/[u]_E$.
\end{lemma}

\begin{lemma}\label{LEMMA_TimidEdgeSeparation}
    Suppose $u,v\in \mathrm{V}(G)$ are such that $u$ is timid and $u\nsim_E v$. Then $u\nsim_{E_\mathrm{t}} v$.
\end{lemma}

\begin{proof}
    Without loss of generality, in view of Lemma \ref{LEMMA_TimidEqvClassQuot}, we may assume that $u\nsim_E u'$ for every $u'\in \mathrm{V}(G)$.

    Striving for a contradiction, suppose that $u\sim_{E_\mathrm{t}}v$. We recursively construct a sequence $\seq{(P_n,F_n):{n\in\omega}}$ such that, for every $n\in\omega$, 
    \begin{itemize}
        \item[(i)] $F_n\subset \mathrm{E}(G)$ is finite;
        \item[(ii)] $\mathcal P_{n+1}$ is a finite set of paths from $u$ to $\mathrm{V}(P_n)$ with edges in $F_n$;
        \item[(iii)] $F_{n}$ separates $u$ from $\mathrm{V}(\mathcal P_n)$; 
        \item[(iv)] $u\sim_{E_\mathrm{t}} w$ in $G\setminus \left(\bigcup_{i< n}F_i\right)$ for some $w\in \mathrm{V}(\mathcal P_n)$. 
    \end{itemize}
    
    For starters, consider a maximal collection $\mathcal P_0$ of pairwise edge-disjoint paths from $u$ to $v$. Since $u\nsim_E v$, $\mathcal P_0$ is finite. 
    
    We claim that there is a $w\in \mathrm{V}(\mathcal P_0)$ such that $u\sim_{E_\mathrm{t}} w$ in $G\setminus \mathrm{E}(\mathcal P_0)$. Suppose not, i.e., that there is a finite $F\subset E_\mathrm{t}$ separating $u$ from $\mathrm{V}(\mathcal P_0)$ in $G\setminus \mathrm{E}(\mathcal P_0)$. Then such $F$ cannot separate $u$ from $\mathrm{V}(\mathcal P_0)$ in $G$, otherwise said $F$ would also separate $u$ from $v$ in $G$, contradicting our hypothesis. Thus, there must be some path $P$ in $G$ connecting $u$ to $\mathrm{V}(\mathcal P_0)$ which does not go through any edge in $F$. Let $w'$ be the first vertex from $\mathrm{V}(\mathcal P_0)$ which appears in $P$. Then the truncation of $P$ from $u$ to $w'$ does not go through any edge in $\mathrm{E}(\mathcal P_0)$, thus $F$ should break $P$, a contradiction. 
    
    Now, If $w\in \mathrm{V}(\mathcal P_0)$ is such that $u\nsim_{E_\mathrm{t}}w$, fix a finite $F_w\subset E_\mathrm{t}$ which separates $u$ from $w$ and then let
    \[
      F_0 = \mathrm{E}(\mathcal P_0)\cup \bigcup_{\substack{w\in \mathrm{V}(\mathcal P_0)\\w\nsim_{E_\mathrm{t}}u}}F_w.
    \]

    
    In this case, because $\bigcup_{\substack{w\in \mathrm{V}(\mathcal P_0)\\w\nsim_{E_\mathrm{t}}u}}F_w\subset E_\mathrm{t}$ is finite, there is still a $w\in \mathrm{V}(\mathcal P_0)$ such that $u\sim_{E_\mathrm{t}} w$ in $G\setminus F_0$. Thus, it is clear that (i)--(iv) hold so far.

    Suppose $(\mathcal P_i,F_i)_{i\le n}$  are defined in a way which satisfies (i)--(iv) thus far. Consider a maximal collection $\mathcal P_{n+1}$ of pairwise edge-disjoint paths from $u$ to $\mathrm{V}(\mathcal P_n)$ in $G\setminus \left(\bigcup_{i\le n}F_i\right)$. Since $u\nsim_E w$ for any $w\in \mathrm{V}(\mathcal P_n)$, $\mathcal P_{n+1}$ is finite. 

    
    We claim that there is a $w\in \mathrm{V}(\mathcal P_n)$ such that $u\sim_{E_\mathrm{t}} w$ in $G\setminus \left(\mathrm{E}(\mathcal P_{n+1})\cup\bigcup_{i\le n}F_i\right)$. Suppose not, i.e., that there is a finite $F\subset E_\mathrm{t}$ separating $u$ from $\mathrm{V}(\mathcal P_{n+1})$ in $G\setminus \left(\mathrm{E}(\mathcal P_{n+1})\cup\bigcup_{i\le n}F_i\right)$. Then such $F$ cannot separate $u$ from $\mathrm{V}(\mathcal P_{n+1})$ in $G$, otherwise said $F$ would also separate $u$ from $v$ in $G$, contradicting our hypothesis. Thus, there must be some path $P$ in $G$ connecting $u$ to $\mathrm{V}(\mathcal P_{n+1})$ which does not go through any edge in $F$. Let $w'$ be the first vertex from $\mathrm{V}(\mathcal P_{n+1})$ which appears in $P$. Then the restriction of $P$ from $u$ to $w'$ does not go through any edge in $\mathrm{E}(\mathcal P_{n+1})$, thus $F$ should break $P$, a contradiction. 
    
    Now, If $w\in \mathrm{V}(\mathcal P_{n+1})$ is such that $u\nsim_{E_\mathrm{t}}w$, fix a finite $F_w\subset E_\mathrm{t}$ which separates $u$ from $w$ and then let
    \[
      F_{n+1} = \mathrm{E}(\mathcal P_{n+1})\cup \bigcup_{\substack{w\in \mathrm{V}(\mathcal P_{n+1})\\w\nsim_{E_\mathrm{t}}u}}F_w.
    \]

    In this case, because $\bigcup_{\substack{w\in \mathrm{V}(\mathcal P_{n+1})\\w\nsim_{E_\mathrm{t}}u}}F_w\subset E_\mathrm{t}$ is finite, there is still a $w\in \mathrm{V}(\mathcal P_{n+1})$ such that $u\sim_{E_\mathrm{t}} w$ in $G\setminus \left(\cup\bigcup_{i\le n+1}F_i\right)$. Thus, it is clear that (i)--(iv) hold so far and our recursion is done. Now consider the graph 
    \[
        H = \bigcup_{n\in\omega}\left(\bigcup\mathcal P_n\setminus \{u\}\right).
    \]
    
    We claim that $H$ is connected. Indeed, note that $v\in \mathrm{V}(H)$. We inductively show that if $w\in \mathrm{V}(\mathcal P_n)$ for some $n\in\omega$, then there is a concatenation of segments of paths in $\bigcup_{k\le n}\mathcal{P}_k$ connecting $w$ to $v$. Suppose that this is true for every $w\in \bigcup_{k\le n} \mathrm{V}(\mathcal P_k)$ and let $w\in \mathcal P_{n+1}$ be given. By construction of $\mathcal{P}_{n+1}$, $w\in \mathrm{V}(P)$ for some path $P$ from $u$ to some $w'\in \mathrm{V}(\mathcal P_n)$. Thus, concatenating the segment of $P$ starting at $w$ with the path $P'$ given by our induction hypothesis to $w'$ we reach the desired conclusion. 
    
    Moreover, $H$ is locally finite: indeed, it follows from (ii) and (iii) that $F_n$ separates $\mathrm{V}(\mathcal P_n)$ from $\mathrm{V}(\mathcal P_{n+1})$ in $H$. The Star-Comb Lemma shows us that $H$ contains some ray $r$ (since $H$ is infinite, by construction) which passes through infinitely many of the $\mathcal P_n$'s. Note that if $w\in \mathrm{V}(\mathcal P_n)\cap r$ and $w'\in \mathrm{V}(\mathcal P_m)\cap r$ for $n\neq m$, then there are $H$-paths $P\in \mathcal P_n$ and $P'\in \mathcal P_m$ passing through $w$ and $w'$, respectively. But, by construction, $\mathrm{E}(\mathcal P_n)\cap \mathrm{E}(\mathcal P_m)=\emptyset$, so it follows that $P$ and $P'$ are edge-disjoint. We at last conclude that $u$ dominates $r$, contradicting our hypothesis that $u$ is timid. 
\end{proof}

\end{document}
